% strat_t3_example — worked example: the dollar-return dollar investor, mid-2022.
% Numbers transcribed from tables/strat_t3_example.pdf (generated by strategy.R).
\footnotesize
\setlength{\tabcolsep}{4pt}
\begin{tabular}{l c c c c l c}
  \toprule
  Month & $\FXGCF^{\mathrm{oos}}_{t}$ &
  $\widehat{\E}_{t}[\rx_{t+12}]$ (\%) & $\widehat{\sigma}_{t}$ (\%) &
  $w_{t}$ & Action & Realised 12m (\%) \\
  \midrule
  2022-05 & $0.17$  & $3.26$ & $6.49$ & $1.22$ & Hold overweight & $-7.9$ \\
  2022-06 & $-0.13$ & $2.29$ & $6.62$ & $0.82$ & Sell ${\sim}32\%$ of bonds $\to$ underweight & $-6.5$ \\
  2022-07 & $-0.65$ & $0.63$ & $6.73$ & $0.22$ & Sell ${\sim}73\%$ of bonds $\to$ near-cash & $-10.4$ \\
  \bottomrule
\end{tabular}
\tabnotes{Each month the investor reads the recursive dollar-return signal
  $\FXGCF^{\mathrm{oos}}_{t}$, updates the real-time forecast
  $\widehat{\E}_{t}[\rx_{t+12}]$ of the next twelve months' dollar excess
  return and the recursive return volatility $\widehat{\sigma}_{t}$, and sets
  the mean--variance target weight,
  $w_{t}=(1/\gamma)\,\widehat{\E}_{t}[\rx_{t+12}]/\widehat{\sigma}^{2}_{t}$;
  ``Action'' is the implied rebalancing. ``Realised 12m'' is the subsequent
  realised twelve-month return, shown only to judge the call ex post---it is
  not used in forming the weight. The deteriorating signal drove a move to
  near-cash ahead of a sharp drawdown.}
